\subsection{What the retrospective transfer analyses establish}
\label{sec:discussion-transfer}
The full-source and nested SPH comparisons use the same predictions. Their large differences therefore concern conditioning on eligibility and the available reference ontology rather than improved fitted weights. In particular, the P4 full-source mapped macro-F1 of 0.0388 and nested value of 0.7278 describe very different statement prevalences. A panel selected to contain comparable numbers of positive and negative references answers whether a score separates those selected groups. It does not estimate the number of false alerts generated when the positive statement is rare in an unselected source. The full-source analysis is useful precisely because it exposes this denominator sensitivity, while remaining limited by the source's statement semantics.

Source-associated reference errors can affect apparent transport independently of waveform shift. A diagnosis may be absent because it was not assessed, because it was encoded differently or because it was genuinely absent. SPH absent-statement comparisons are retained as that declared mapping; LUDB unknown incidental statements are not silently treated as negatives. PTB Diagnostic provides explicit healthy controls for named infarct territories, but leaves many ECGs unknown for each territory. These different reference structures explain why one unified all-database macro-F1 would be misleading. Reporting known reference states and source-specific endpoints is more informative than forcing a common score onto incompatible labels.

These completed analyses remain retrospective. The model has already been inspected on SPH and several public resources, and some development experiments fit new heads or policies. Source novelty alone does not restore an untouched evaluation after such inspection. The independent external study remains necessary to quantify performance under a prospectively frozen analytical contract. Its eventual results will require their own cohort flow, overlap accounting, uncertainty analysis and limitations, rather than being appended as another row in the current selected-development table.

\subsection{Interpretation of specialization and uncertainty}
\label{sec:discussion-specialization}
Specialist branches provide a clear place to define output vocabularies and operating points. They can also make it easier to describe which endpoints were directly supervised and which were merely exposed by a global head. This organizational advantage does not establish superior diagnostic performance over a matched single-head baseline. The literature tables describe related methods and evaluation units; there was no controlled common-cohort comparison against all cited systems. Parameter counts establish the scale of the executed release architecture, not accuracy per parameter, energy efficiency or deployment suitability.

The two fitted backbones also limit the meaning of a shared representation. The rhythm and pathology packages implement the same backbone design, but the released V3F and V3AG files contain separately fitted parameters. They do not represent one joint forward pass with a common embedding for every reported result. A future deployment comparison should measure peak memory, startup, waveform preparation and complete latency for the actual combined route. The small LUDB fallback timing excludes model startup and the cached primary-stage runtime, so it cannot establish prospective throughput or a service-level response guarantee.

Confidence-based withholding should be evaluated with the same care as classification. The SPH policy reduced conditional known-endpoint error among accepted ECGs, while withholding correct opportunities and positive references. Its larger full-denominator loss at the selected illustrative withholding cost shows why accepted accuracy alone cannot establish benefit. Different clinical costs might favor different policies, but those costs would require an explicitly justified utility model and prospective evaluation. The current score-distance policy does not incorporate emergency urgency, patient symptoms, repeat testing or clinician availability.

The executed LUDB fallback demonstrates a narrower engineering fact: a withheld ECG can be rescored through the P6 waveform path and receive some additional membership calls. It leaves NSR, PAC and all P4 endpoints unresolved for those ECGs. Reporting complete coverage as 100\% would hide that limitation. The correct descriptions are complete P8/P4 coverage of 152/200 and at least partial output on 200/200, with 77 of the original 924 known opportunities unresolved. No evidence from that execution establishes referral effectiveness or clinical safety.

\subsection{Reproducibility at distinct levels}
\label{sec:reproducibility}
The available evidence supports several different verification tasks. Archived predictions permit independent recalculation of selected metrics and figures. Published checkpoint lengths and hashes identify the released files. Instrumented executions establish tensor shapes and repeatability for a reconstructed configuration. Surviving run summaries describe selected optimization and training activity. None of these tasks alone establishes the complete sequence of data, software, intermediate states and decisions that produced the original models.

The strict P4 historical comparison is an example of this distinction. Matching all 962 threshold decisions did not make 836 out-of-tolerance probabilities numerically equivalent. Threshold agreement is useful for describing that panel's retained calls, but it does not authenticate the historical probability-calibration calculation or guarantee agreement near a decision boundary in other records. Likewise, a correct SHA-256 digest authenticates a byte sequence; it cannot establish that a particular scheduler job read those bytes or that its training patients were disjoint from later evaluation patients.

For future releases, a stronger chain would bind each training job to a versioned source tree, resolved configuration, dependency versions, patient/signal-component rosters, input checkpoint digests, optimization state, output digests and evaluation ledger. A protected record-level ledger can coexist with publicly shared aggregates and code, subject to source licenses. This would make a later distinction between numerical replay, retraining reproducibility and independent scientific replication easier to assess. The present revision reports the verification level actually available and preserves unresolved lineage as an explicit limitation.

Separate corrective and autonomous development records were inspected to determine whether any experimental candidate replaced the original release. Their recorded decisions retained the incumbent and did not authorize promotion. These studies involve different candidate constructions, preprocessing or repeatedly inspected selection banks. Their optimization scores are therefore not combined with the selected P6/P4 estimates or presented as independent external results. The five original released artifacts remain the model identities of the primary analysis.
